\scanpage{196}

Nous allons déterminer la structure multiplicative de $A(G)$. Les \unsure{hom.} \diagram{\begin{tikzcd}[ampersand replacement=\&] F\arrow[r,dash,"i"] \& G\arrow[r,dash,"j"] \& H \end{tikzcd}} définissent des homomorphismes

\diagram{\begin{tikzcd}[ampersand replacement=\&] A(H)\arrow[r,"j^*"] \& A(G)\arrow[r,"i^*"] \& A(F) \end{tikzcd}}

$A(H)$ est l'algèbre \inserted{sur $Z$} engendrée par les $\sigma_n$ $(n\in Z)$ soumis aux relations $\sigma_n\sigma_{n'}=\sigma_{n\unsure{-|+}n'},\ \sigma_0=1$. Il est plongé dans $A(G)$ par $j^*$, par

\[ j^*(\sigma_n)=e_{n\,1}. \]

$A(F)$ est \deleted{\unsure{l'algèbre en}} une algèbre ayant une base formée d'éléments $e_n$ $(n\geq1)$ \deleted{\unsure{et}} $e_1=1$. D'ailleurs $A(F)=Z[e_2]$. On a

\[ i^*(e_{nk})=e_k \]

Identifions $\Lambda=A(H)$ à une sous-algèbre de $A(G)$. Posons $\xi=e_{0\,2}$. On se propose de prouver que

\[ A(G)=\Lambda[\xi]\qquad\text{(alg. des polynômes \unsure{en $\xi$,|à une $\xi$,} \unsure{cff.|coeff.} \unsure{dans $\Lambda$)|de $\Lambda$)}} \]

i.e. que les $\xi^q$ $(q=0,1,\ldots)$ forment une \uline{base} de $A(G)$ sur $\Lambda$. \deleted{\unsure{Pour}}\inserted{Comme} les $e_{0\,1},e_{0\,2},e_{0\deleted{\unsure{j}}3}\ldots$ forment une telle base, il \deleted{\unsure{on}}, il suffit d'établir qu'on a des formules

\[ \xi^q=\sum_{i=1}^{q+1}c_{qi}e_{0\,i}\qquad(c_{qi}\in\Lambda,\ c_{q,q+1}=1)\tag*{(1)} \]

% archive sheet number: 195
